\section{Bewertung der Nutzerzufriedenheit}
\label{ch:Diskussion:Zufriedenheit}

Auf Grundlage der jahrezehntelangen Erfahrung mit dem alten System ist anhand der Befragung festzustellen, dass die Nutzer einer Modernisierung ihrer Routine teilweise sehr kritisch gegenüber stehen. Auch wenn das neue System sehr intuitiv sein sollte, folgt daraus nicht unbedingt, dass der Nutzer es mag.

Mithilfe der Befragung wurde eine Meinung verzeichnet, wie geeignet die befragten Personen das Dash\-board für den gestellten Anwendungsfall halten.

Allgemein sind die Nutzer sehr zufrieden mit dem neuen Dash\-board. 80\% der Fragen haben bezüglich des Dash\-boards ein signifikant besseres Ergebnis erhalten. Nach Ausfilterung der zwei Ausnahmen, welche die mit Abstand größte Expertise aufwiesen (s. \autoref{ch:Diskussion:Effektivität}), steigt dieser Wert auf 100\%.

Abgeleitet von den vorangegangenen Bewertungen gibt es keine Diskrepanz zwischen den objektiven Testergebnissen und der subjektiven Nutzererfahrung.

Da die Stichprobe in Alter, Berufserfahrung und bis auf wenige Ausnahmen der Expertise sehr homogen ausfällt, kann keine weitere Ableitung auf eine der moderierenden Variablen getätigt werden. 

% Allgemein
%Nutzer sind zufrieden damit\\
%--> Es lässt sich zusammenfassend sagen, dass die Nutzer sehr zufrieden mit dem neuen Dashboard sind.\\
%--> Prümper\\
%--> Gibt es eine Diskrepanz zwischen Nutzen und Nutzbarkeit eines System (obj / subj)?\\
%\textbf{Decken sich die Ergebnisse der Umfrage und des Tests?}\\
%-	wie sich gemessene Werte und Umfrageergebnisse decken / bestätigen\\
%-	Bewerten in jeder Nuance\\
%  --> moderierende Variablen, Effizienz/Zeit, Effektivität/Genauigkeit, Umfrage\\
%  --> alles in die Kategorien der ISO-Norm???\\

Zum Abschluss des Versuchs wurden die Nutzer gefragt, ob sie sich vorstellen können, das Dash\-board zukünftig zu nutzen (s. \autoref{fig:DB_Nutzung:boxplot}). Bis auf die selben zwei Personen, die ein klares 'nein' antworteten, haben alle weiteren Testpersonen mit 'eher ja' und 'ja' geantwortet. Auch der t-Test bestätigt eine statistische Signifikanz dieser Aussage.
